\section*{四、实验探讨与思考题}

\subsection*{1. 实验心得}

通过本次实验，我学会了用牛顿环和劈尖装置测量平凸透镜曲率半径和薄片厚度，也了解了其干涉条纹的成因。在实验过程中，应使显微镜移动方向准确通过牛顿环圆心，并保持十字叉丝与暗环边缘相切的判断标准一致，以减小定位误差。观测过程中，为避免读数显微镜空程差的影响，应单向移动；目测左侧时选择暗圆外沿，右侧时选择暗圆内沿。此外，应尽量避免触碰牛顿环和劈尖装置，并保持固定螺钉松紧适当，以免空气膜厚度和条纹形状发生变化。

\subsection*{2. 思考题}

\begin{enumerate}[itemsep=0.2em, topsep=0.3em]
  \item \textbf{如何消除读数显微镜引起的空程差？}测量时应使手轮始终沿同一方向转动，避免中途反向，从而消除螺杆与螺母间隙对读数的影响。
  \item \textbf{为什么随着牛顿环半径增大，条纹越来越密？}牛顿环的空气膜厚度随半径增大而加快变化，因此相邻两级条纹所对应的半径差逐渐减小，表现为越靠外侧条纹越密。
  \item \textbf{在透射光中是否也可以观察到干涉条纹？}可以。透射光中同样能观察到干涉条纹，其明暗与反射光互补，即反射暗环对应透射亮环；但透射条纹的对比度通常较低。
\end{enumerate}
\setstretch{1.45}
